\chapter{染色体把遗传规律搬进细胞}

\begin{kaichang}
找一张你家的全家福，或者翻翻手机里一家人的合影。先看一眼最明显的事：你像爸爸，也像妈妈——也许是爸爸的眼睛形状配上妈妈的脸型。再看看你的兄弟姐妹（如果没有，就看看你和堂表兄弟姐妹的合影）：你们明明来自同一对父母，却没有人长得一模一样。一个双眼皮，一个单眼皮；一个自然卷，一个直发；连酒窝都可能只长在一个人脸上。

现在请你猜：同一对父母，用同一套“配方”，为什么每次做出的“产品”都不一样？是完全随机抓阄，还是有一套可以计算的规则？第 63 章里，孟德尔用豌豆告诉你：规则确实存在，而且可以写成比例。但孟德尔本人从头到尾没见过他所说的“遗传因子”长什么样、藏在细胞的什么地方。这一章，我们要把那些看不见摸不着的“因子”，搬进一个真实的、能在显微镜下看到的住所——染色体。
\end{kaichang}

\section{孟德尔的“因子”住在哪里}

先把第 63 章的结论压缩成三句话：性状由成对的遗传因子（今天我们叫等位基因）决定；形成生殖细胞时，成对的因子彼此分开，各进各的（分离定律）；不同性状的因子通常各自独立地组合（独立分配定律）。孟德尔靠种豌豆和数豆子推出了这些，比例漂亮得像数学定理。

但请你站在 1900 年前后一位生物学家的立场上想一想，这套理论有一个让人睡不着的漏洞：\textbf{因子在哪里}？分离定律说，成对的因子在形成生殖细胞时会“分开”。这是个非常具体的物理动作——分开，意味着它们本来待在一起，然后被某种力量分到了两边。什么东西在细胞里成对存在，又恰好在形成生殖细胞时分开？如果找不到这样东西，孟德尔的因子就永远只是一个数学幽灵：比例是对的，但没人知道是什么东西在执行这些比例。

这一章的故事，就是把这个幽灵抓进显微镜里。而抓它的线索，来自一个看似毫不相干的领域：细胞分裂。

\section{把镜头推进细胞核}

第 57 章参观细胞时，你见过细胞核——细胞里那个被双层膜包起来的“档案室”。档案室里存放的，是细胞全部遗传信息的载体：DNA。不过这一章我们先不急着看 DNA 的分子结构（那是第 65、66 章的主角），先看它的\textbf{打包方式}。

一个普通人体细胞里的 DNA，如果全部拉直首尾相接，长度大约 2 米——比你本人还高。而这 2 米长的分子，要塞进直径只有约 \ud{6}{\mu m}（百万分之六米）的细胞核里。做个类比：相当于把一根 200 千米长的细线，整整齐齐收进一颗乒乓球。随便乱塞必然打结，打结的线没法读、没法抄。细胞的解决方案是：把 DNA 一圈一圈缠在一种叫\textbf{组蛋白}的蛋白质线轴上，再把缠好的线轴层层折叠、捆扎。DNA 加蛋白质的这整套打包物，平时松散地铺开，像一碗面条，叫\textbf{染色质}；到细胞要分裂时，它进一步浓缩成一条条短粗、显微镜下清晰可见的“小棒”，这才是大名鼎鼎的\textbf{染色体}。所以请记住：染色质和染色体不是两种东西，是同一份材料在不同时段的两种打包状态——就像同一床棉被，平时摊开，搬家时捆紧。

你的每个体细胞（除了成熟红细胞等少数例外）里有 46 条染色体，分成 23 对。每一对里，一条来自父亲的精子，一条来自母亲的卵子。这一对叫\textbf{同源染色体}：它们大小、形状相同，相同位置上装着控制同一类性状的基因——比如都在某个位置有一段与眼睛色素有关的基因。但注意，“同一类性状”不等于“同一个版本”：父亲给的那条上可能是“双眼皮版本”，母亲给的那条上可能是“单眼皮版本”。同一位置的不同版本，就是第 63 章说的等位基因。23 对里，22 对男女都一样，叫常染色体；第 23 对是性染色体，女性是两条 X，男性是一条 X 加一条小得多的 Y。

最后一个名词，也是最容易搞混的一个。细胞准备分裂前，会先把每条染色体完整复制一份。复制之后，两份一模一样的拷贝并不马上分开，而是像拉链没拉开一样，在中部一个叫\textbf{着丝粒}的位置连在一起。这两份连在一起的拷贝，叫\textbf{姐妹染色单体}。请把这三个概念的关系记牢：同源染色体是“一条来自爸、一条来自妈”的一对；姐妹染色单体是“同一条染色体复印出来的双份”。一个是“配对”，一个是“复印”，完全不同。后面你会看到，减数分裂的全部精妙，就藏在这两对概念的先后登场里。

\section{有丝分裂：一丝不苟的抄写与平分}

你的身体从一个受精卵长到大约 37 万亿个细胞，靠的是有丝分裂。第 59 章讲过细胞怎样合作形成身体，这里我们把镜头对准“分裂”这个动作本身，因为它的逻辑简单得漂亮。

有丝分裂要解决的问题是：一个细胞变成两个，每个新细胞都必须拿到一套完整无缺的遗传档案。细胞的解法分两步。第一步，\textbf{先抄后分}：分裂之前，把 46 条染色体全部复制一遍——每条变成由着丝粒相连的一对姐妹染色单体。此刻细胞里有 46 条“双线”染色体，信息总量翻了一倍。第二步，\textbf{对半平分}：细胞里会架起一套由蛋白质纤维组成的“牵引系统”（纺锤体），从两端拉住每条染色体的着丝粒，把连体的一对姐妹染色单体拆散，一份拉去左边，一份拉去右边。然后细胞从中间缢裂成两个。

结果是什么？左边 46 条，右边 46 条，每条都来自原来那对姐妹中的一份——也就是说，两个子细胞拿到的档案\textbf{逐字相同}。你的皮肤细胞分裂产生皮肤细胞，肝细胞产生肝细胞，你伤口边缘的细胞分裂补上缺口，靠的都是这套“复制一份、均分两份”的机制。有丝分裂维持的不是单个细胞，而是一条绵延几十年的细胞谱系：今天你肝细胞里的遗传档案，和你三岁时肝细胞里的，是同一套档案的一次次抄写件（抄写偶尔也会出错，那是第 72 章的话题）。

请特别留意：有丝分裂里，同源染色体从头到尾各过各的——来自爸的那条和来自妈的那条，从不配对，也不交换任何东西。平分的是姐妹染色单体，不是同源染色体。记住这一点，下一节全靠它。

\section{减数分裂：减半、洗牌与交换}

如果生殖细胞也靠有丝分裂产生，麻烦就大了：精子和卵子各有 46 条染色体，受精后孩子就有 92 条，孙子 184 条……代代翻倍，几代人之后细胞核就装不下了。染色体数目必须有个“减半”的机制，受精时两套“半份”拼回一整套。这个机制就是\textbf{减数分裂}，它只在产生精子和卵子的细胞里发生。

减数分裂的口诀是：\textbf{复制一次，分裂两次}。第一次分裂（减数第一次分裂）分的是同源染色体——来自爸的和来自妈的彼此分开，各进一个子细胞；第二次分裂（减数第二次分裂）才分姐妹染色单体，和有丝分裂类似。一减再减，最终每个配子（精子或卵子）只含 23 条染色体：每对同源染色体中的一条。受精时，精子的 23 条遇上卵子的 23 条，拼回 46 条——不多不少，代代稳定。

现在你明白孟德尔的分离定律是谁在执行了：所谓“成对的因子彼此分开”，就是减数第一次分裂时同源染色体的分离。基因就坐在染色体上，染色体分开，基因自然分开。幽灵有了身体。

更精彩的是独立分配定律。减数第一次分裂的中期，23 对同源染色体在细胞中央排队，每对各自决定“爸的那条朝哪边、妈的那条朝哪边”——而各对之间的朝向互不相关，就像 23 枚硬币各自独立抛掷。你 1 号染色体分到“爸版”，完全不影响 2 号分到哪一版。这就是独立分配定律在细胞里的样子：不同染色体上的基因，确实各自独立地洗牌。

\begin{qiaoliang}
来算一笔账：只考虑 23 对染色体的随机朝向，一个人能产生多少种染色体组合不同的配子？每对同源染色体有 2 种选法，23 对就是 $2^{23}$ 种，约等于 838 万种。也就是说，你父亲能产生约 838 万种精子，你母亲能产生约 838 万种卵子，你和你亲兄弟姐妹恰好抽到同一套组合的概率，比中彩票头奖低得多。而且这还没算下一节要讲的交换——算上交换，可能的组合数实际上远超任何人口数量。同一对父母，永远做不出两个一样的孩子（同卵双胞胎除外，他们来自同一个受精卵的分裂）。全家福上的“人人不同”，不是意外，是数学的必然。
\end{qiaoliang}

减数分裂还有最后一道工序，也是最聪明的一道：\textbf{交换}。在减数第一次分裂的早期，配对在一起的同源染色体紧紧相贴，然后——来自爸的那条和来自妈的那条，会在对应的位置断开，交换一段，再精确接上。这不是断裂事故，而是细胞动用专门的酶完成的精确操作，交换前后，每条染色体上的基因总数一个不多、一个不少，只是“爸版”和“妈版”的片段被拼到了同一条上。交换的结果：进入配子的那条染色体，既不是纯爸版也不是纯妈版，而是一条全新的混编版。洗牌的层次又多了一层：不只整染色体在洗牌，染色体内部的片段也在洗牌。

\begin{shiyan}
\textbf{看一看正在分裂的细胞（需要显微镜）}。材料：学校或商售的“洋葱根尖有丝分裂”永久装片一片、光学显微镜一台。步骤：先用低倍镜找到根尖生长点区域——那里细胞小而密、近似正方形；再转高倍镜，慢慢扫视，寻找染色体形态清晰的细胞。观察点：你会看到有的细胞里染色体排成一条中线（即将平分），有的已经分成两团向两端移动，还有的核膜完整、染色质松散（暂时没分裂）。数一数：大约多少比例的细胞正处在分裂状态？安全提示：永久装片已固定染色，无化学风险，但玻片边缘可能锋利，拿取时捏侧边；不可在家用菜刀切洋葱根尖自行染色——常用的染色剂有腐蚀性和毒性，想自己做只能在学校实验室由老师带领。这个观察的意义在于：你亲眼看到的是几十亿年来每一个多细胞生物身体里每时每刻都在发生的事。
\end{shiyan}

\section{给染色体标上字母}

到这一节，符号终于可以登场了。第 63 章你用 A、a 表示等位基因，现在我们知道它们住在哪：每个基因占据染色体上一个确定的位置，叫\textbf{座位}。你可以把染色体想成一条长街，基因是街上的门牌号。来自爸的街和来自妈的街，门牌号一一对应（同源），但门牌里住的“版本”可以不同。

减数分裂时发生的事，用符号写就是：一对同源染色体，一条标 A，一条标 a，彼此分离，各进一个配子——配子要么带 A，要么带 a，概率各半。两对位于不同染色体上的基因 A/a 和 B/b，则随着两对染色体各自独立朝向，配子得到 AB、Ab、aB、ab 四种，比例均等。这和第 63 章的棋盘格完美对上，但现在每个符号背后都有了一根实实在在的“小棒”。

可如果两个基因住在\textbf{同一条}染色体上呢？假设 A 和 B 在同一条街上，父亲给的那条是“AB”，母亲给的那条是“ab”。这时它们不再独立分配：没有交换时，这条街是整条传下去的，配子只能得到 AB 或 ab，得不到 Ab 和 aB。这种现象叫\textbf{连锁}——两个基因被“拴”在同一根小棒上，倾向于捆绑遗传。

但交换会部分解开这根绳：交换发生在两个座位之间时，就能拼出 Ab 和 aB 这样的新组合。两个座位离得越远，它们之间发生交换的机会越大，新组合就越多；离得越近，越难被拆开。换句话说，\textbf{两个基因被拆开的频率，暴露了它们在染色体上的距离}。这句话听起来不起眼，它却是人类画出第一批基因地图的全部原理——下面的证据故事会告诉你这是怎么发生的。

\begin{wuqu}
“染色体交换是细胞出了故障，把染色体弄断了。”——恰恰相反。交换是减数分裂的\textbf{标准工序}：每次形成精子或卵子，每对同源染色体通常都会发生一到数次交换，由专门的酶切割、对接、修复，位置对应、分毫不差。没有交换反而要出事——交换除了制造新组合，还帮助同源染色体在分离前稳固地“牵手”定位，交换太少与不分离（后面讲的非整倍体）风险升高有关。把交换当成故障，就像把“洗牌”当成“撕牌”：洗牌当然不是撕牌，它是让牌局公平且常新的关键步骤。真正的事故是另一回事：染色体断裂后接错位置（易位）、整段丢失或重复——那些属于结构变异，和第 72 章的突变主题相关。
\end{wuqu}

\section{证据：从蝗虫到一只白眼果蝇}

\begin{zhengju}
1902 到 1903 年，两条线索在显微镜下汇合了。美国人萨顿观察蝗虫的生殖细胞，德国人鲍维里用海胆卵做实验，两人各自注意到同一件事：染色体在减数分裂中的行为，和孟德尔因子的行为，严丝合缝地平行——因子成对，染色体成对；因子分离，染色体分离；因子独立组合，染色体独立朝向；受精时因子恢复成对，染色体也恢复成对。他们提出一个大胆猜想：\textbf{基因就在染色体上}。这就是“染色体遗传学说”。但请注意，此时它只是一个猜想：行为平行不等于同住一处。影子和人动作一致，也可能只是巧合。

把猜想钉死的是摩尔根和他的果蝇。1910 年，摩尔根在培养的红眼果蝇里发现了一只白眼雄蝇。他让这只雄蝇和红眼雌蝇交配，第一代全是红眼——符合孟德尔，红色显性。让第一代互相交配，第二代红眼与白眼的比例约 3:1——还是符合孟德尔。但有一个细节怎么也对不上：\textbf{白眼个体几乎全是雄蝇}。如果控制眼色的基因在某条常染色体上，性别和眼色不该有任何瓜葛。摩尔根提出了唯一的自洽解释：这个基因在 X 染色体上。雄蝇只有一条 X，来自母亲，上面带什么版本就表现什么；雌蝇有两条 X，要两条都带白眼版本才会是白眼——所以白眼在雄蝇里常见得多。一个基因，第一次被“定位”到了一条具体的染色体上。猜想变成了实物。

故事还有下半场。摩尔根的学生斯特蒂文特注意到：同在 X 染色体上的几个基因，捆绑遗传的紧密程度各不相同——有的几乎从不分开，有的经常因交换而分开。他想到上一节那个不起眼的原理：拆开频率就是距离。1913 年，还是大学生的斯特蒂文特把六个基因的交换频率换算成间距，画出了历史上第一张基因连锁图——一张“染色体街道地图”，上面只有相对位置和距离，却准确预言了后来直接观测到的基因顺序。这套方法的伟大之处在于：它用繁殖计数这种“笨办法”，测出了当时任何显微镜都看不见的分子间距。当然，今天的基因组测序早已把每条街的门牌号抄录得清清楚楚，但第一张地图，是用纸笔和果蝇画出来的。
\end{zhengju}

\section{这套机器也会出错：非整倍体}

任何机械流程都有出错的时候，减数分裂也不例外。最重要的一类错误叫\textbf{不分离}：减数分裂时，某一对同源染色体（或姐妹染色单体）该分开却没分开，结果一个配子拿到 24 条染色体——某号染色体多了一条，另一个配子则少了一条。多拿一条的配子受精后，胚胎就会有 47 条染色体，某号染色体有三份而不是两份。染色体数目不是整倍数，这叫\textbf{非整倍体}。

你多半听说过它的一个实例：21 号染色体多一条，即 21 三体，也就是唐氏综合征。它是新生儿中最常见的染色体数目异常之一，会带来程度不一的发育差异和健康风险。为什么是 21 号最常见？一个可能的原因是：其他染色体的三体对发育的干扰更大，多数在很早的孕期就无法继续，只有最小的几条染色体（21 号是最小的常染色体之一）多一条时，胚胎还能发育到出生。这不是“21 号容易出错”，而是“出错后还能出生的，大多轮到 21 号”——观察到的频率，是出错率和存活率两道筛子共同筛过的结果。

统计上还观察到一个现象：母亲年龄越大，生育 21 三体孩子的概率越高。这与卵细胞的特殊经历有关：女性的卵细胞在母亲还是胎儿时就已启动减数第一次分裂，然后\textbf{暂停}在那里几十年，直到排卵时才继续完成。染色体的“牵手”结构要维持几十年不出差错，难度可想而知。这里要强调两点：这是概率层面的群体规律，不是对任何个人的断言；而且受精需要精子与卵子双方，染色体不分离也可能发生在精子形成过程中。如果你或家人遇到相关的医学疑问，请咨询专业医生，科普书不提供个体诊疗建议。

非整倍体给我们的真正启示是：遗传信息不仅要求“内容正确”，还要求“份数正确”。基因产物像菜谱里的配料，多一勺少一勺都会改变整个系统的平衡——剂量本身就是信息的一部分。

\section{再看那张全家福}

现在回到开场的问题：同一对父母，为什么做不出两个一样的孩子？

拆开看，你已经有了完整的链条。父亲产生精子时，23 对同源染色体随机分半，$2^{23}$ 种组合；母亲产生卵子时同样如此。受精是两次“大抽奖”的相遇，仅染色体整条的随机分配，就给出约 $2^{23} \times 2^{23}$，也就是七十万亿种可能的组合。再叠加上每条染色体内部的交换——父亲传给孩子的那条 1 号染色体，其实是爷爷奶奶两条 1 号混编而成的新版本——组合数实际上大到没有名字。

所以全家福上的每个人都是独一无二的：你从父亲那里拿到的，是他从他父母那里抽来、又经过交换重洗的一副牌；你的兄弟姐妹拿到的是同一副牌的另一次洗牌结果。你们共享大约一半的等位基因——所以相像；但没有任何两个人抽到完全相同的一手——所以不同。像与不像，不是命运也不是巧合，是减数分裂这台“洗牌机”的出厂设置。

顺便回应一个你可能听过的问题：“为什么孩子更像爸还是更像妈，总被亲戚争论？”从染色体层面看，答案平淡无奇：除了性染色体和线粒体那点例外，来自父母双方的常染色体份额严格相等。至于“看起来更像谁”，取决于哪些基因恰好是显性、哪些性状更显眼、以及观察者先看鼻子还是先看眼睛——那更多是注意力的游戏，不是遗传学的偏心。

\section{染色体里，到底谁带着信息}

这一章我们把孟德尔的因子安顿进了染色体，故事看似圆满。但如果你是个爱抬杠的读者，此刻应该提出一个尖锐的问题：染色体是\textbf{两种}东西打包在一起的——DNA 和蛋白质。到底哪一种才是真正的遗传物质？

别急着说 DNA。站在 1930 年的知识水平，押注蛋白质的人理由相当充分：蛋白质由 20 种氨基酸组成，花样繁多，折叠出的形状千变万化，看起来完全有资格储存复杂信息；而 DNA 当时被认为只是四种核苷酸的单调重复，像个不起眼的结构材料。把遗传信息的重任交给谁，不是一个能靠猜解决的问题——它需要实验，而且需要那种能把两种分子干净分开、逐个检验的漂亮实验。

谁做到了？怎样做到？过程中有哪些争论、误判，还有一位女科学家被低估的贡献？下一章——第 65 章，我们来讲这场科学史上著名的“身份鉴定案”：是谁证明 DNA 携带遗传信息。

\begin{tiaozhan}
斯特蒂文特的基因地图为什么能画成直线？关键在于：当两个基因距离不远时，它们之间发生交换的频率近似正比于距离——距离翻倍，交换机会大致翻倍。于是可以定义一个“图距”单位：两个座位间每 100 个后代中平均出现 1 个交换产物，图距就是 1 个单位（今天叫 1 厘摩，cM，纪念摩尔根）。假设 A、B 间交换率 5\%，B、C 间交换率 3\%，且 B 在中间，那么 A、C 间的交换率约是多少？第一直觉是 $5\%+3\%=8\%$，但实际会略低——因为若在 A-B 和 B-C 之间\textbf{各}发生一次交换（双交换），A、C 两端的组合恰好被换回原型，统计时漏掉了。双交换概率约 $0.05 \times 0.03 = 0.15\%$，所以观测值约 $8\% - 2\times0.15\% = 7.7\%$。这正是遗传学家能发现“双交换”现象、并反过来确认基因直线排列的方法。距离更远时比例关系会明显失真，要靠更精细的函数修正。这段跳过不影响主线，但如果你将来看到任何一张“基因图谱”，你会知道上面的每一格距离，都是成千上万次繁殖计数换来的。
\end{tiaozhan}

\section*{练一练}

\begin{sikaoti}
\item 画一对同源染色体：一条来自父方、标为 A（显性），一条来自母方、标为 a。先画出复制后带姐妹染色单体的样子（着丝粒要标出来），再画出减数第一次分裂和第二次分裂后，四个配子各得到什么。在图旁写一句话：姐妹染色单体分离和同源染色体分离，分别对应哪一步？
\item 有丝分裂和减数分裂都包含“姐妹染色单体分离”这一步，但一个叫“维持”、一个叫“减半”。用染色体数目的变化（从 46 条出发）写清两种分裂的起点和终点，并解释：为什么减数分裂必须先发生一次“同源染色体分离”，姐妹单体的分离才有减半意义？
\item 一对夫妻的第一个孩子像爸爸的单眼皮，第二个孩子像妈妈的双眼皮。亲戚说：“第三个孩子该轮到像爸爸了。”用减数分裂的随机分配说明这个推理错在哪里；如果该性状由一对等位基因控制，且母亲为杂合，第三个孩子为单眼皮的概率是多少？
\item 基因 A 与 B 连锁在同一条染色体上，父亲那条是 AB，母亲那条是 ab。不考虑交换时，他们的孩子能从父亲精子中得到的组合有哪几种？考虑交换后又多了哪几种？如果 A、B 相距很远，交换频繁，后代的组合比例会越来越接近什么经典比例？这说明连锁和独立分配之间是什么关系？
\item 画一张“信息流图”：从“一对同源染色体”出发，经过减数分裂、受精，到“孩子的 46 条染色体”，用箭头标出信息在哪些环节被减半、在哪些环节被重新组合、在哪个环节恢复成双份。
\item 21 三体的患者细胞里有 47 条染色体。请写出导致这一结果的两种可能的不分离情形（分别发生在减数第一次分裂和第二次分裂），并说明：为什么显微镜下数染色体条数，无法区分错误发生在哪一次分裂？
\end{sikaoti}
